% 2026, Prof. Dr. Matthias Jung (DL9MJ)
% DIPOL
\begin{tikzpicture}

% Hauptachse
\begin{axis}[
    name=mainaxis,
    width=\linewidth,
    height=0.66*\linewidth,
    xmode=log,
    ymode=log,
    xmin=0.01,
    xmax=100,
    ymin=0.01,
    ymax=1000000,
    domain=0.01:100,
    samples=500,
    clip=false,
    xlabel={
        Distanz, normiert:
        $\displaystyle
        x=\frac{d}{\lambda/(2\pi)}
        $ 
    },
    ylabel={
        elektrische Feldstärke, normiert:
        $\lvert E\rvert/E_0$
    },
    tick label style={
        font=\normalsize
    },
    xtick={0.01,0.1,1,10,100},
    xticklabels={
        {$0{,}01$},
        {$0{,}1$},
        {$1$},
        {$10$},
        {$100$}
    },
    grid=both,
    %axis on top,
    legend style={
        at={(0.98,0.98)},
        anchor=north east,
        fill=white,
        draw=black,
        font=\large
    }
]

% Elektrisches Feld:
%
% |E|/E0 =
% sqrt[(1/x^3 - 1/x)^2 + 1/x^4]
%
\addplot[
    DARCorange,
    very thick
]
{
    sqrt(
        (1/x^3 - 1/x)^2
        + 1/x^4
    )
};
% Magnetisches Feld:
%
% |H|/(E0/Z0) =
% sqrt[1/x^4 + 1/x^2]
%
\addplot[
    DARCgreen,
    very thick
]
{
    sqrt(
        1/x^4
        + 1/x^2
    )
};
% Markierung bei r = lambda/(2 pi), also x=1
\draw[
    black,
    dashed,
    line width=1pt
]
(axis cs:1,0.01)
--
(axis cs:1,1000000);

\node[
    anchor=north west,
]
at (axis cs:1.05,200000)
{
    $\displaystyle d=\frac{\lambda}{2\pi}$
};

% Erklärung Nahfeld
\node[
    anchor=west,
    fill=white,
    draw=DARCorange,
    align=center
]
at (axis cs:0.060,70000)
{
    $\displaystyle |E|\sim\frac{1}{d^3}$
};
\node[
    anchor=west,
    fill=white,
    draw=DARCgreen,
    align=center
]
at (axis cs:0.035,2)
{
    $\displaystyle |H|\sim\frac{1}{d^2}$
};


% Grenze Fernfeld: r = 4 lambda
\draw[
    black,
    dashed,
    line width=1pt
]
(axis cs:{8*pi},0.01)
--
(axis cs:{8*pi},1000000);

\node[
    anchor=north west,
]
at (axis cs:{8*pi*1.03},200000)
{
    $\displaystyle d=4\lambda$
};

% Erklärung Fernfeld
\node[
    anchor=center,
    fill=white,
    draw=black,
    align=center
]
at (axis cs:25,1)
{
    $\displaystyle |E|,|H|\sim\frac{1}{d}$
};

% Bereiche
\node[
    anchor=south,
    font=\scriptsize\bfseries
]
at (axis cs:0.1,1300000)
{Reaktives Nahfeld};

\node[
    anchor=south,
    font=\scriptsize\bfseries
]
at (axis cs:5,1300000)
{Strahlendes Nahfeld};

\node[
    anchor=south,
    font=\scriptsize\bfseries
]
at (axis cs:50,1300000)
{Fernfeld};

\end{axis}

% Rechte y-Achse für das normierte H-Feld
\begin{axis}[
    width=\linewidth,
    height=0.66*\linewidth,
    at={(mainaxis.south west)},
    anchor=south west,
    xmode=log,
    ymode=log,
    xmin=0.01,
    xmax=100,
    ymin=0.01,
    ymax=1000000,
    axis x line=none,
    axis y line*=right,
    ylabel={
        magnetische Feldstärke, normiert:
        $\lvert H\rvert/(E_0/Z_0)$
    },
    tick label style={
        font=\normalsize
    },
    grid=none
]
\end{axis}
\end{tikzpicture}